On construit une pile cuivre-zinc en reliant, par un pont salin, un bécher
contenant $\qty{20}{\mL}$ de la solution de sulfate de cuivre(II) à
$\qty{0,20}{\M}$ dans laquelle plonge une lame de cuivre à un bécher contenant
$\qty{20}{\mL}$ de la solution de sulfate de zinc(II) de même concentration dans
laquelle plonge une lame de zinc.
\[
    \ce{Cu^{2+}_{(aq)} + Zn_{(s)} <=> Cu_{(s)} + Zn^{2+}_{(aq)}},
    \qquad \mathrm{K}=10^{37}
\]

On réalise un circuit série comprenant cette pile, un conducteur ohmique de
résistance $\qty{10}{\ohm}$ et un milliampèremètre.
\begin{enumerate}
    \item Quel est le rôle du pont salin~?
    \item En appliquant le critère d'évolution, déterminer la polarité des
          électrodes.
    \item Donner les demi-équations des réactions aux électrodes ainsi que
          l'équation de la réaction associée à la transformation qui a lieu dans
          la pile.
\end{enumerate}

